\chapter{Entrelazador con tensores simétricos sin traza y reducción transversal \texorpdfstring{\(12\to5\to5\to2\)}{12→5→5→2}}
\label{chap:parte-iv-64}

El marco icosaédrico construye una isometría equivariante hacia el espacio de tensores simétricos sin traza. La proyección transversal realiza \(12\to5\to5\to2\), conserva el portador pentacomponente y localiza el subespacio radiativo de dimensión \(2\).

\section{Entrelazador con el espacio de tensores sim\'etricos sin traza}

Sea \(\mathcal V\subset S^2\) el conjunto de los doce v\'ertices unitarios
del icosaedro. Para cada \(v\in\mathcal V\), def\'inase
\begin{equation}
\label{rm:eq:gravity-Qv}
Q_v=vv^T-\frac13I_3\in
\operatorname{STF}_3:=\operatorname{Sym}^2_0(\R^3).
\end{equation}
El mapa de marco
\begin{equation}
\label{rm:eq:gravity-Gamma0}
\Gamma_0\!\left(\sum_{v\in\mathcal V}x_ve_v\right)
=\sum_{v\in\mathcal V}x_vQ_v
\end{equation}
es \(A_5\)-equivariante y satisface
\begin{equation}
\label{rm:eq:gravity-tight-frame}
\Gamma_0\Gamma_0^*=\frac85I_{\operatorname{STF}_3},
\qquad
\Gamma_0=\Gamma_0P_5.
\end{equation}

\begin{theorem}[Entrelazador tensorial normalizado]
\label{rm:thm:gravity-gamma5}
La aplicaci\'on
\begin{equation}
\label{rm:eq:gravity-gamma5}
\boxed{
\Gamma_5:=\sqrt{\frac58}\,\Gamma_0|_{V_5}:
V_5\xrightarrow{\ \simeq\ }\operatorname{STF}_3}
\end{equation}
es una isometr\'ia \(A_5\)-equivariante. Su adjunto es
\begin{equation}
\label{rm:eq:gravity-gamma5-adjoint}
\Gamma_5^*Q=\sqrt{\frac58}
\sum_{v\in\mathcal V}(v^TQv)e_v.
\end{equation}
\end{theorem}

\begin{proof}
La equivarianza sigue de \(Q_{gv}=gQ_vg^T\). El operador
\(\Gamma_0\Gamma_0^*\) conmuta con la representaci\'on irreducible de
\(A_5\) sobre \(\operatorname{STF}_3\), por lo que es escalar. Su traza
es
\[
\sum_{v\in\mathcal V}\|Q_v\|_F^2
=12\cdot\frac23=8;
\]
al dividir por la dimensi\'on cinco se obtiene \(8/5\). De aqu\'i,
\(\Gamma_5\Gamma_5^*=I\). Como dominio y codominio tienen dimensi\'on
cinco, tambi\'en \(\Gamma_5^*\Gamma_5=I\). La f\'ormula del adjunto
procede de \(\langle Q,Q_v\rangle_F=v^TQv\).
\end{proof}

La normalizaci\'on mediante el marco icosa\'edrico elimina una coincidencia
puramente cardinal: construye el mapa que transporta el sumando finito al
portador continuo de esp\'in dos.

\section{Proyecci\'on transversal sin traza}

Para \(n\in S^2\), sea \(\Pi(n)=I-nn^T\). Definimos
\begin{equation}
\label{rm:eq:gravity-lambda-tt}
\boxed{
\Lambda(n)Q
=\Pi Q\Pi-\frac12\Tr(\Pi Q\Pi)\Pi.}
\end{equation}

\begin{theorem}[Reducci\'on de cinco componentes a dos]
\label{rm:thm:gravity-tt}
La restricci\'on de \(\Lambda(n)\) a \(\operatorname{STF}_3\) es el
proyector ortogonal sobre
\begin{equation}
\label{rm:eq:gravity-htt}
\mathcal H_{\rm TT}(n)
=\{Q\in\operatorname{STF}_3:Qn=0\}.
\end{equation}
En particular,
\begin{equation}
\label{rm:eq:gravity-lambda-properties}
\Lambda(n)^2=\Lambda(n)=\Lambda(n)^*,
\qquad \rank\Lambda(n)=2.
\end{equation}
Para \(n=e_3\), la imagen est\'a generada por
\begin{equation}
\label{rm:eq:gravity-plus-cross}
E_+=\frac1{\sqrt2}\operatorname{diag}(1,-1,0),
\qquad
E_\times=\frac1{\sqrt2}
\begin{pmatrix}0&1&0\\1&0&0\\0&0&0\end{pmatrix}.
\end{equation}
\end{theorem}

\begin{proof}
Como \(\Pi n=0\), la salida es transversal. Su traza vale
\[
\Tr(\Pi Q\Pi)-\frac12\Tr(\Pi)\Tr(\Pi Q\Pi)=0,
\]
pues \(\Tr\Pi=2\). Si \(Qn=0\) y \(\Tr Q=0\), entonces
\(\Pi Q\Pi=Q\) y \(\Tr(\Pi Q\Pi)=0\), por lo que
\(\Lambda(n)Q=Q\). La f\'ormula es autoadjunta para el producto de
Frobenius. En el marco \(n=e_3\) su matriz en una base STF adaptada es
\(\operatorname{diag}(1,1,0,0,0)\), lo que prueba el rango.
\end{proof}

La cadena completa queda, por tanto,
\begin{equation}
\label{rm:eq:gravity-12-5-5-2}
\boxed{
V_{12}\xrightarrow{P_5}V_5
\xrightarrow{\Gamma_5}\operatorname{STF}_3
\xrightarrow{\Lambda(n)}\mathcal H_{\rm TT}(n),
\qquad 12\longrightarrow5\longrightarrow5\longrightarrow2.}
\end{equation}
El acoplamiento reducido
\begin{equation}
\label{rm:eq:gravity-effective-map}
\mathbb G_{{\rm TT},a}(n)
=G_a(\theta_{108}^{\rm circ})\Lambda(n)\Gamma_5P_5
\end{equation}
tiene rango dos.


\paragraph{Descomposición helicoidal del espectro y separación de los modos masivo y no masivos.} La reducción permite distinguir el espectro helicoidal completo de una realización masiva y sus dos helicidades sin masa.
